\subsection{Reference-controlled completion in a fixed vertical range}

Put
\[
 c_R=
 \begin{cases}
 3,&R=\mathbb Z,\\
 2p-1,&R=\mathbb Z_{(p)}.
 \end{cases}
 \tag{AT.181}
\]
The first positive cohomology of a free strict PD
identity disk over $R$ occurs no earlier than
degree $c_R+1$. This follows from the actual
arity differential $R\xrightarrow{n}R$ for an
odd-to-even disk, as verified below.

\begin{theorem}[AT27: Proved; actual completion controlled by a common reference and paths]
\label{AT27}
Let $2\leq d\leq c_R$, $S=\Gamma_R(T)$, and let
$A_i$ be a specified finite nonempty family of
torsionfree admissible theories over $R$.
Let $X,D$ be strict simplicial PD cochain
$S$-algebras with torsionfree terms, coefficient
$H^0$ in every simplex and weight, and rational
coefficient comparison in every simplex and weight.
Suppose $D$ has integral coefficient comparison
through form $d$ for weights at least $L$, uniformly
in the simplex degree.
Specify the actual maps
\[
 g:X\longrightarrow D,\qquad r_i:D\longrightarrow A_i,
 \qquad F_i:X\longrightarrow A_i^{\Delta[1]},
 \qquad e_0F_i=r_i g .
 \tag{AT.182}
\]
Write $f_i=e_1F_i$; these are the maps to be retained.
Set
\[
 W=\max\bigl\{1,d+1,L,\lambda_{A_i}(d+3),
                              \mu_{A_i}(d+3):i\bigr\}.
 \tag{AT.183}
\]
There are actual inclusions $X\hookrightarrow E$ and
$D\hookrightarrow D'$, strict maps
\[
 g':E\longrightarrow D',\qquad r_i':D'\longrightarrow A_i,
 \qquad F_i':E\longrightarrow A_i^{\Delta[1]},
 \qquad e_0F_i'=r_i'g',
 \tag{AT.184}
\]
extending all of (AT.182), with the following properties.
\begin{enumerate}
\item $E,D'$ have torsionfree terms, coefficient
      $H^0$ in every weight and simplex, and rational
      coefficient comparison in every weight and simplex.
\item Both have integral coefficient comparison
      through form $d$ for $w\geq W$, uniformly
      in the simplex degree.
\item Every horizontal row is eventually acyclic.
      One may take the row cutoff in form $q$ to be
      $\max\{q+1,\mu_{A_i}(q):i\}$ and then take
      maxima through a requested form range.
\item The endpoint maps $e_1F_i'$ are actual
      extensions of the originally prescribed $f_i$.
\end{enumerate}
If the support condition $q\leq w$ holds initially,
it holds in these constructions.
No assertion of eventual vertical exactness above
degree $d$ is made.
No identification of the low-form truncation of
$E$ with that of $X$ is asserted: the construction
includes explicit new generators in those forms.
\end{theorem}

\begin{proof}
We use the initial-endpoint ideal
\[
                 K_i=\ker(e_0:A_i^{\Delta[1]}\to A_i).
 \tag{AT.185}
\]
Its rows are contractible in every weight and
form degree by the order-preserving minimum
contraction in A-002.
The same proposition, with its displayed
finite-range hypercohomology argument, gives
\[
 H^q((K_i)_n^{w,*})=0
 \quad(0\leq q\leq d,\ w\geq W,\ n\geq0).
 \tag{AT.186}
\]
Indeed its proof in degree $d$ uses the vertical
and row cutoffs through $d+3$, all included in
(AT.183). These are uniform in $n$.
The path terms and $K_i$ are torsionfree:
they are finite-limit submodules of the
corresponding torsionfree simplex groups.
Flat rationalization commutes with these
finite diagrams, so (AT.186) also gives
the needed rational vanishing of $H^0(K_i)$.

First consider a current source still denoted
$X$, with the actual maps (AT.182), and a
homogeneous torsion pair in an arbitrary simplex:
\[
 z\in X_n^{w,q},\qquad dz=0,\qquad
 b\in X_n^{w,q-1},\qquad db=mz,
 \quad 2\leq q\leq d,\quad w\geq W,\quad m\ne0.
 \tag{AT.187}
\]
No normalization or unproved homogeneous
normalization choice is needed.
The representable free simplicial module
$R[\Delta[n]]$ has
$\operatorname{Hom}(R[\Delta[n]],M)=M_n$.
Thus (AT.187) is an actual pair of simplicial
parameter maps on that representable.

Choose elements in the reference with
\[
 du_D=g(z),\qquad dv_D=mu_D-g(b).
 \tag{AT.188}
\]
They exist in $D_n$ by its comparison through
$d$, since $q,q-1$ are positive.
Let $c:A_i\to A_i^{\Delta[1]}$ denote constant paths.
The difference
$F_i(z)-d\,c r_i(u_D)$ is a closed element
of $(K_i)_n^{w,q}$.
By (AT.186) choose $h_{u,i}$ in degree $q-1$
whose differential is that difference, and put
\[
                      U_i=c r_i(u_D)+h_{u,i}.
 \]
Then $dU_i=F_i(z)$ and $e_0U_i=r_i(u_D)$.
Similarly
$mU_i-F_i(b)-d\,c r_i(v_D)$ is closed in
$(K_i)_n^{w,q-1}$.
Choose its primitive $h_{v,i}$ and put
\[
                      V_i=c r_i(v_D)+h_{v,i}.
 \tag{AT.189}
\]
It satisfies $dV_i=mU_i-F_i(b)$ and
$e_0V_i=r_i(v_D)$.

Adjoin free $R[\Delta[n]]$ copies of generators
$u,v$, of forms $q-1,q-2$ and weight $w$, with
\[
                           du=z,\qquad dv=mu-b.
 \tag{AT.190}
\]
The assignments $g(u)=u_D,g(v)=v_D$ and
$F_i(u)=U_i,F_i(v)=V_i$ define actual strict
maps on this entire simplicial free extension.
They satisfy the endpoint equation exactly.
All faces and degeneracies of the chosen values
are the values of the corresponding representable
parameters. This is why no separate compatibility
choices among simplex degrees occur.

Rationally the actual triangular change
$u'=u-b/m$ gives $du'=0,dv=mu'$.
In each simplex this is a free rational
graded-commutative identity disk, with its
explicit monomial contraction.
For a set of such pairs the finite-variable
contractions and their filtered unions give
the same rational coefficient comparison.
Take the image of the extension in its
rationalization. The old source remains
injective, and all the indicated maps factor
because their targets are torsionfree.
All new weights are at least $W$.
Every new closed zero-form is rationally a
coefficient; its image in $D$ is an integral
closed zero-form, and hence is the same
integral coefficient.
This proves coefficient $H^0$ throughout.

We now give the saturation used here, which is
controlled by the actual reference map rather
than by a choice of a simplex-zero presentation.
During a saturation stage the rational source
algebra is fixed.
For a current subalgebra $Z$ in it, adjoin
in each weight $w\geq W$ every rational
zero-form $a$ such that
\[
                   da\in Z,\qquad g_{\mathbb Q}(a)\in D.
 \tag{AT.191}
\]
The conditions are stable under addition,
coefficient multiplication, faces and
degeneracies. Form the PD subalgebra generated
by these roots and iterate this operation
countably many times.
The differential preserves each generated
algebra because the derivative of every
new generator was already integral.

Every existing map $g$ extends by restricting
its rationalization. We check the path maps
with their actual endpoint values.
For a root $a$, the element
\[
 \xi_i=F_{i,\mathbb Q}(a)
                     -c\,r_i(g_{\mathbb Q}(a))
 \tag{AT.192}
\]
lies in $(K_i)_{n,\mathbb Q}^{w,0}$.
Its differential is integral and belongs
to $(K_i)_n^{w,1}$.
By (AT.186) it has an integral primitive there.
The difference from $\xi_i$ is a rational
closed zero-form in $K_i$, hence zero.
Thus $\xi_i$ itself is integral, and so is
$F_{i,\mathbb Q}(a)$.
Products and divided powers remain integral.
All relations and the endpoint equations are
retained because these extensions are
restrictions of the rational algebra maps.

This saturation kills every closed one-form
in the stated weights. Let $t$ be such a form
in one stage. Rational exactness gives $da=t$.
In $D$ choose an integral primitive $b_D$ of
$g(t)$, using $H^1(D)=0$ in this weight.
Then $g_{\mathbb Q}(a)-b_D$ is a rational
coefficient. Subtract that coefficient from
$a$. The resulting primitive satisfies
(AT.191), so it is adjoined at the next stage.
Any closed zero-form in the saturation is
rationally a coefficient whose image in
$D$ is integral, so coefficient $H^0$ is
retained. Rationalization is unchanged;
the new terms are torsionfree.
No hypothesis $X_0=S$ or claim that its
first skeleton stays fixed is used.

Alternate this saturation with the simultaneous
extensions (AT.190) for all nonzero cohomology
classes in forms two through $d$, every
simplex $n$ and every weight $w\geq W$.
All these classes have nonzero annihilators
because the rational coefficient comparison
holds. Choose their actual representatives
and relations (AT.187).
Each stage uses the same reference $D$ and
the same bound (AT.183).
Take the compatible countable union.
Every positive closed form in the specified
range appears in one stage and receives a
primitive in a later one; degree one uses
(AT.191). This gives an actual extension
$X\hookrightarrow E_0$ with all the desired
maps and vertical comparison through $d$.

It remains to restore the horizontal rows
without losing this finite vertical range.
We give the simultaneous construction on
the source and reference, retaining the paths.
Let $z\in N_n(E_0^{w,q})$ be a horizontal
cycle with
\[
                  w\geq M(q):=
                    \max\{q+1,\mu_{A_i}(q):i\}.
 \tag{AT.193}
\]
Adjoin to the reference a C21 free identity
disk whose normalized top generator $a_D$
has form $q$, degree $n+1$, weight $w$ and
boundary $g(z)$.
The latter is an actual normalized cycle.
Choose, independently in each target, an
actual normalized horizontal primitive
$a_i\in N_{n+1}(A_i^{w,q})$ of $r_i g(z)$.
These choices extend $r_i$ to that new
reference disk, including its differential.

Adjoin the corresponding free disk to the
source, with top $a$ and boundary $z$, and
send it to $a_D$ under $g$.
For the path value, the difference
\[
                   F_i(z)-\partial c(a_i)
                     \in N_n(K_i^{w,q})
 \tag{AT.194}
\]
is a horizontal cycle.
The rows of $K_i$ are contractible in every
weight, so choose
$h_i\in N_{n+1}(K_i^{w,q})$ with this boundary.
Assign
$F_i(a)=c(a_i)+h_i$ and assign its derivative
to the differential partner.
These are actual free-disk maps, with
$e_0F_i(a)=a_i=r_i g(a)$.

Also adjoin reference disks for any other
normalized horizontal cycles in the reference
in the same range (AT.193).
Their maps to the targets extend by the same
horizontal exactness. Existing source maps
and paths then compose with reference inclusion.
Iterate these two operations on all source
and reference cycles.
This is the C21 free-disk construction with
the additional explicitly given maps
(AT.194). Every required cycle is killed in
a later stage, so both final algebras have
the row cutoffs asserted in the statement.

We verify preservation of the vertical range
in this last construction.
In every simplex a relative free C21 cell
adjoins free identity-disk complements.
An even-to-odd disk has the integral pairs
$\gamma_n(u)\mapsto\gamma_{n-1}(u)du$.
For an odd generator of form $r\geq1$ the
arity-$n$ pair is
\[
 R\,(u\gamma_{n-1}(du))
       \xrightarrow{\,n\,}R\,\gamma_n(du)
 \quad\text{in forms }n(r+1)-1,\ n(r+1).
 \tag{AT.195}
\]
Over $\mathbb Z$ its first nonzero positive
cohomology is no earlier than form four.
Over $\mathbb Z_{(p)}$ the integers $1,\ldots,p-1$
are units, so that degree is at least $2p$.
Here is a direct tensor check, valid also
when the current terms are merely torsionfree.
Decompose one disk into its arity-zero copy
of $R$ and the displayed arity pairs.
Every unit-$n$ pair is contractible after
tensoring with the current complex, using
its explicit two-term inverse homotopy.
Every nonunit pair has its lower form
$a=n(r+1)-1$ at least $c_R$.
Its tensor with a nonnegative current
complex has no term in degrees below $a$.
In degree $a=c_R$, a cocycle would be a
closed current zero-form killed by $n$:
its two differential components are
$dx$ and $nx$.
Thus this group is
$\ker(n:H^0(\text{current})\to H^0(\text{current}))=0$,
because the coefficient $H^0=S$ is weightwise
free over $R$. Pairs beginning later do not
contribute in the required range at all.
Consequently tensoring one disk leaves
cohomology through $c_R$ unchanged, including
its coefficient $H^0$.
Iterate this argument over any finite
set of free disk generators and then take
their filtered union.
The decomposition respects weights, so
it proves the claim weightwise without
assuming $S$-flatness or free current terms.

These free extensions also preserve
torsionfreeness, coefficient $H^0$ and
rational coefficient comparison by their
pointwise disk descriptions.
Condition (AT.193) gives the support bound
for each new differential partner.
The earlier vertical cells and zero-form
roots retain that bound by (AT.183).
Taking the final compatible unions proves
all assertions (AT.184) and completes the
construction.
\end{proof}

Applied to the actual maps in AT26, this gives
an all-simplex local completion through degree
three, with the original endpoint maps retained.
The reference in that application is the actual
algebra (AT.174)--(AT.175), not an inferred
sequence of separate simplexwise choices.

The finite-form gluing requirement C34 in
SA08 requires a strict isomorphism
$(P_p)_{<d_p}\cong(E_p)_{<d_p}$ preserving its
specified target maps. The present construction
instead supplies the displayed inclusions and
can adjoin generators in those very low forms.
It has not supplied that C34 isomorphism.
